\section*{Bab \ref{ch:b2:realfun} --- Fungsi Peubah Real}

\begin{solution}{exo:b2:realfun:1}
$\lfloor x\rfloor$: lompatan berukuran $1$ di setiap bilangan bulat. $x -
\lfloor x\rfloor$: lompatan berukuran $-1$ di bilangan bulat (dengan
limit kiri $1$ dan nilai $0$). $\lfloor x\rfloor + \sqrt{x - \lfloor
x\rfloor}$: di bilangan bulat $n$, limit kirinya $(n - 1) + 1 = n$ dan
nilainya $n$, jadi \emph{\omterm{def:b2:metric:continuity}{kontinu}} di mana-mana (sebab akar kuadratnya
memperbaiki lompatannya), walaupun tak dapat diturunkan di bilangan
bulat. Adapun fungsi berlompatan rasional tadi: setelah konstruksinya
dibatasi pada pencacahan bilangan diadik, ia melompat sebesar $2^{-n}$
persis di bilangan rasional diadik ke-$n$ dan \omterm{def:b2:metric:continuity}{kontinu} di tempat lain.
\end{solution}

\begin{solution}{exo:b2:realfun:2}
Andaikan $f$ yang naik punya ketakkontinuan di titik dalam $a$: maka
$f(a^-) < f(a^+)$ (\cref{thm:b2:realfun:monotone}) dan peta
$I$ melewatkan interval \omterm{def:b2:metric:topology}{terbuka} tak kosong
$\intoo{f(a^-)}{f(a^+)}$ kecuali mungkin satu nilai $f(a)$:
jadi peta sembarang subinterval yang memuat $a$ di bagian dalamnya bukan
sebuah interval (sebab ia berjurang pada sedikitnya satu sisi $f(a)$).
Hal ini bertentangan dengan sifat nilai antaranya. Adapun ketakkontinuan
di ujungnya disingkirkan dengan cara sama lewat jurang sepihak.
\end{solution}

\begin{solution}{exo:b2:realfun:3}
$x\ln x$: turunan keduanya $\frac1x > 0$, jadi cembung. $\ln(1 +
\eu^x)$: turunannya $\frac{\eu^x}{1 + \eu^x} = 1 - \frac{1}{1 +
\eu^x}$, yang naik, jadi cembung. $\sqrt{1 + x^2}$: turunan keduanya
$(1 + x^2)^{-3/2} > 0$, jadi cembung. $x^3$: tidak cembung pada $\R$
(sebab $f'' = 6x$ berganti tanda); ia cembung hanya pada $\R_+$.
\end{solution}

\begin{solution}{exo:b2:realfun:4}
Andaikan $f(a) \neq f(b)$, katakanlah $f(b) > f(a)$ dengan $a < b$
(kasus $f(b) < f(a)$ setangkup, dengan memandang ke kiri). Untuk $x > b$,
ketaksamaan kemiringan (\cref{lem:b2:realfun:slopes}) pada $a < b < x$ memberi
\[
\frac{f(x) - f(a)}{x - a} \geq \frac{f(b) - f(a)}{b - a} = m > 0
\quad\Longrightarrow\quad
f(x) \geq f(a) + m(x - a) \xrightarrow[x\to+\infty]{} +\infty,
\]
yang bertentangan dengan keterbatasannya di atas. Jadi $f$ konstan.

Tentang asimtotnya: jika $f(x) - (\alpha x + \beta) \to 0$ di $+\infty$
dan $f(x) - (\alpha' x + \beta') \to 0$ di $-\infty$, maka fungsi
cembung $g(x) = f(x) - (\alpha x + \beta)$ terbatas di atas di dekat
$+\infty$; lalu kecembungan beserta asimtot di $-\infty$ (yang memaksa
$\alpha' \leq \alpha$ lalu $\alpha' = \alpha$ dengan membandingkan
kemiringannya di $\mp\infty$: sebab kemiringan fungsi cembung naik)
membuat $g$ terbatas di atas pada seluruh $\R$, jadi konstan $= 0$ pada
limitnya: sehingga $f$ afin.
\end{solution}

\begin{solution}{exo:b2:realfun:5}
Fungsi $f'$ monoton, jadi menurut \cref{thm:b2:realfun:monotone}
satu-satunya ketakkontinuan yang mungkin berupa lompatan, dengan limit
sepihak yang ada di mana-mana. Menurut \cref{cor:b2:realfun:nojumps},
sebuah turunan tak punya ketakkontinuan berupa lompatan. Karenanya $f'$
sama sekali tak punya ketakkontinuan: jadi \omterm{def:b2:metric:continuity}{kontinu}.
\end{solution}

\begin{solution}{exo:b2:realfun:6}
Ambil logaritmanya:
\[
\frac1p \ln\Bigl(\sum_i \lambda_i x_i^p\Bigr)
= \frac1p \ln\Bigl(1 + p\sum_i \lambda_i \ln x_i +
O(p^2)\Bigr)
= \sum_i \lambda_i \ln x_i + O(p)
\xrightarrow[p \to 0^+]{} \sum_i \lambda_i \ln x_i ,
\]
dengan memakai $\sum\lambda_i = 1$ dan $\ln(1 + u) = u + O(u^2)$.
Setelah dieksponensialkan diperoleh rata-rata geometrinya. Kini untuk
setiap $p \in \intoo{0}{1}$, ketaksamaan rata-rata pangkat
(\cref{ex:b2:realfun:powermeans}, dengan eksponen $p < 1$) memberi
\[
\Bigl(\sum_i \lambda_i x_i^p\Bigr)^{1/p} \leq \sum_i\lambda_i x_i ;
\]
lalu melewatkan $p \to 0^+$ di ruas kirinya menghasilkan $\prod_i
x_i^{\lambda_i} \leq \sum_i \lambda_i x_i$: itulah ketaksamaan AM--GM berbobot.
\end{solution}

\begin{solution}{exo:b2:realfun:7}
Dengan $\varphi(t) = t\ln t$ (yang cembung, sebab $\varphi'' = \frac1t >
0$) dan bobot $q_i$ di titik $t_i = \frac{p_i}{q_i}$:
\[
\sum_i p_i \ln\frac{p_i}{q_i}
= \sum_i q_i\, \varphi\Bigl(\frac{p_i}{q_i}\Bigr)
\;\geq\; \varphi\Bigl(\sum_i q_i \frac{p_i}{q_i}\Bigr)
= \varphi(1) = 0 ,
\]
lewat Jensen (\cref{thm:b2:realfun:convexreg}~(3)). Kesamaan pada
Jensen bagi fungsi cembung \emph{tegas} memaksa semua titik
$t_i$ berimpit: jadi $\frac{p_i}{q_i}$ konstan, dan setelah dijumlahkan,
konstantanya $1$: yakni $p = q$. (Besaran ini --- yaitu divergensi
Kullback--Leibler --- kembali pada dunia \cref{ch:b2:randomvar}.)
\end{solution}

\begin{solution}{exo:b2:realfun:8}
\emph{Bobot diadik.} Secara induktif pada $m$: kasus $m = 1$ tak lain
hipotesisnya. Untuk bobot $\lambda = \frac{k}{2^{m+1}}$ (dengan $k$
ganjil), tulis $\lambda = \frac12(\lambda_1 + \lambda_2)$ dengan
$\lambda_j = \frac{k \mp 1}{2^{m+1}}$, yang keduanya berpenyebut $2^m$
setelah disederhanakan; maka
\[
f\bigl(\lambda x + (1{-}\lambda)y\bigr)
= f\Bigl(\tfrac{u + v}{2}\Bigr)
\leq \frac{f(u) + f(v)}{2}
\leq \lambda f(x) + (1 - \lambda) f(y),
\]
dengan $u = \lambda_1 x + (1 - \lambda_1)y$ dan $v = \lambda_2 x +
(1-\lambda_2)y$, memakai kecembungan titik-tengah lalu hipotesis
induksinya pada $u, v$.

\emph{Pelewatan ke limit.} Untuk $\lambda \in \intcc{0}{1}$ yang
sembarang, ambil bilangan diadik $\lambda_n \to \lambda$: kekontinuan
$f$ beserta kekontinuan pemetaan afinnya melewatkan ketaksamaan
$f(\lambda_n x + (1-\lambda_n)y) \leq \lambda_n f(x) +
(1-\lambda_n)f(y)$ ke limitnya: jadi $f$ cembung.
\end{solution}

\begin{solution}{exo:b2:realfun:9}
Untuk $0 < x < y$: ketaksamaan kemiringan
(\cref{lem:b2:realfun:slopes}) di titik $0 < x < y$ memberi
\[
\frac{f(x) - f(0)}{x} \leq \frac{f(y) - f(0)}{y},
\qquad\text{yakni}\qquad
\frac{f(x)}{x} \leq \frac{f(y)}{y} +
f(0)\Bigl(\frac1x - \frac1y\Bigr).
\]
Karena $f(0) \leq 0$ dan $\frac1x - \frac1y > 0$, suku terakhirnya
$\leq 0$: jadi $\frac{f(x)}{x} \leq \frac{f(y)}{y}$. Karenanya $x
\mapsto \frac{f(x)}x$ naik.

Kesuperaditifan untuk $f(0) = 0$: untuk $x, y > 0$ (sebab kasus dengan
peubah nol bersifat sepele),
\[
f(x) = x\,\frac{f(x)}{x} \leq x\,\frac{f(x+y)}{x+y},
\qquad
f(y) \leq y\,\frac{f(x+y)}{x+y},
\]
lewat kemonotonan yang baru dibuktikan tadi; lalu menjumlahkannya
memberi $f(x) + f(y) \leq f(x+y)$.
\end{solution}

\begin{solution}{exo:b2:realfun:10}
Untuk $x, y \in I$ dan $\lambda \in \intcc01$: kecembungan $f$,
lalu kemonotonan $g$, lalu kecembungan $g$ memberi
\[
g\bigl(f(\lambda x + (1{-}\lambda)y)\bigr)
\leq g\bigl(\lambda f(x) + (1{-}\lambda)f(y)\bigr)
\leq \lambda\,g(f(x)) + (1{-}\lambda)\,g(f(y)).
\]
Contoh penyangkal tanpa kemonotonan: $g(t) = -t$ bersifat cembung
(sebab afin) tetapi turun, sedangkan $f(x) = x^2$ cembung, dan $g \circ
f = -x^2$ cekung tegas.
\end{solution}

\begin{solution}{exo:b2:realfun:11}
\emph{Ketaksamaan kirinya:} misalkan $m = \frac{a+b}2$ lalu ambil garis
penyangga di $m$ (\cref{thm:b2:realfun:convexreg}~(2)): $f(t) \geq
f(m) + \mu(t - m)$ untuk setiap $t \in \intcc ab$. Setelah diintegralkan
atas $\intcc{a}{b}$: suku linearnya berintegral $\mu\int_a^b(t -
m)\dd t = 0$ (berkat kesetangkupan di sekitar $m$), jadi $\int_a^b f
\geq (b - a)f(m)$.

\emph{Ketaksamaan kanannya:} pada $\intcc ab$, kecembungannya membatasi
$f$ oleh tali busurnya: $f(t) \leq f(a) + \frac{f(b) - f(a)}{b - a}(t -
a)$. Setelah diintegralkan: $\int_a^b f \leq (b-a)f(a) + \frac{f(b) -
f(a)}{b - a}\cdot\frac{(b-a)^2}2 = (b - a)\,\frac{f(a) + f(b)}2$. Lalu
bagilah dengan $b - a$.
\end{solution}

\begin{solution}{exo:b2:realfun:12}
Misalkan $a - \delta < a \leq x < y \leq b < b + \delta$, semuanya di $I$.
Dua kali penerapan ketaksamaan kemiringan
(\cref{lem:b2:realfun:slopes}), mula-mula pada $a - \delta < a \leq x <
y$, lalu pada $x < y \leq b < b + \delta$, memberi
\[
\frac{f(a) - f(a - \delta)}{\delta}
\leq \frac{f(y) - f(x)}{y - x}
\leq \frac{f(b + \delta) - f(b)}{\delta}
\]
(sebab kemiringan tali busurnya naik ketika kedua ujungnya bergerak ke
kanan). Karenanya tiap kemiringan tali busur di dalam $\intcc ab$
terjebak di antara dua bilangan tetap, dan
\[
\abs{f(y) - f(x)} \leq K\,\abs{y - x},
\qquad
K = \max\Bigl(\Bigl|\frac{f(a) - f(a-\delta)}{\delta}\Bigr|,
\Bigl|\frac{f(b+\delta) - f(b)}{\delta}\Bigr|\Bigr):
\]
jadi $f$ bersifat \omterm{def:b2:metric:continuity}{Lipschitz} pada $\intcc ab$. Setiap titik $I$ yang
\omterm{def:b2:metric:topology}{terbuka} punya ruas bermarjin semacam itu di sekitarnya: jadi \omterm{def:b2:metric:continuity}{Lipschitz}
lokal, sehingga (sekali lagi) \omterm{def:b2:metric:continuity}{kontinu} pada $I$.
\end{solution}

\begin{solution}{pb:b2:realfun:1}
\textbf{1.} \emph{Cembung $\Rightarrow$ $f'$ naik:} untuk $a <
b$, ketaksamaan kemiringan memberi, untuk $h > 0$ yang kecil,
$\frac{f(a+h) - f(a)}h \leq \frac{f(b) - f(a)}{b-a} \leq \frac{f(b) -
f(b-h)}{h}$; lalu melewatkan $h \to 0$ memberi $f'(a) \leq \frac{f(b) -
f(a)}{b - a} \leq f'(b)$. \emph{Sebaliknya}, jika $f'$ naik
dan $x < y < z$: teorema nilai rata-rata memberi $c_1 \in
\intoo{x}{y}$ dan $c_2 \in \intoo yz$ dengan
\[
\frac{f(y) - f(x)}{y - x} = f'(c_1) \leq f'(c_2)
= \frac{f(z) - f(y)}{z - y},
\]
lalu ketaksamaan kemiringan tiga titik itu, yang diterapkan dengan $y =
\lambda x + (1 - \lambda)z$, tertata ulang menjadi ketaksamaan
kecembungannya. Untuk kasus $C^2$: $f'' \geq 0$ bila dan hanya bila $f'$ naik.

\textbf{2.} Jika $f'' > 0$, maka $f'$ naik tegas, dan
perhitungan nilai rata-rata di atas memberi ketaksamaan \emph{tegas}
antara kedua kemiringan tali busurnya: jadi cembung tegas.
\emph{Penyangga tegas:} di titik dalam $a$ dengan kemiringan penyangga
$m$, seandainya $f(x_0) = f(a) + m(x_0 - a)$ untuk suatu $x_0 \neq a$,
maka pada ruas dari $a$ ke $x_0$ garis penyangga dan tali busurnya
berimpit, dan kecembungan tegas di titik tengahnya memberi
$f\bigl(\frac{a + x_0}2\bigr) < f(a) + m\,\frac{x_0 - a}2$, yang
bertentangan dengan ketaksamaan penyangganya. Jadi $f(x) > f(a) + m(x -
a)$ untuk setiap $x \neq a$. \emph{Jensen yang tegas:} dengan $a =
\sum\lambda_ix_i$, merata-ratakan ketaksamaan penyangganya memberi
$\sum\lambda_if(x_i) \geq f(a)$, dengan kesamaan bila dan hanya bila tiap
sukunya menjadi kesamaan, yakni bila dan hanya bila setiap $x_i = a$.

\textbf{3.} $(-\ln)'' = \frac1{t^2} > 0$; $(t^r)'' = r(r -
1)t^{r-2}$, yang positif untuk $r > 1$ dan negatif untuk $0 < r < 1$;
serta $\exp'' = \exp > 0$. Semuanya tegas menurut pertanyaan 2.

\textbf{4.} Kasus $ab = 0$ bersifat sepele. Untuk $a, b > 0$,
kecekungan $\ln$ di titik $a^p, b^q$ dengan bobot
$\frac1p, \frac1q$ memberi
\[
\ln\Bigl(\frac{a^p}p + \frac{b^q}q\Bigr) \geq
\frac1p\ln(a^p) + \frac1q\ln(b^q) = \ln(ab),
\]
dan $\ln$ naik: jadi $ab \leq \frac{a^p}p + \frac{b^q}q$.
Kesamaannya bila dan hanya bila kedua titiknya berimpit (lewat
kecekungan tegas), yakni $a^p = b^q$.

\textbf{5.} Kecekungan $\ln$ dengan bobot $\lambda_i$ memberi
$\ln\bigl(\sum\lambda_ix_i\bigr) \geq \sum\lambda_i\ln x_i =
\ln\prod x_i^{\lambda_i}$; lalu eksponensialkan. Kesamaannya bila dan
hanya bila semua $x_i$ sama (pertanyaan 2). Jalur pada
\cref{exo:b2:realfun:6} memperoleh ketaksamaan yang sama sebagai limit
rata-rata pangkat; sedangkan di sini ia satu penerapan Jensen --- jadi
kotak perkakasnya memang punya kelebihan pasokan.

\textbf{6.} Jika $a = 0$ atau $b = 0$, ketaksamaannya sepele.
Normalkan: dengan mengganti $a$ oleh $a/\norm a_p$ dan $b$ oleh $b/\norm
b_q$, kita boleh mengandaikan $\norm a_p = \norm b_q = 1$ lalu wajib
menunjukkan $\sum\abs{a_ib_i} \leq 1$. Young suku demi suku memberi
\[
\sum_i\abs{a_i}\abs{b_i} \leq
\sum_i\Bigl(\frac{\abs{a_i}^p}{p} +
\frac{\abs{b_i}^q}{q}\Bigr) = \frac1p + \frac1q = 1 .
\]
Kesamaannya bila dan hanya bila tiap ketaksamaan Young-nya ketat: yakni
$\abs{a_i}^p = \abs{b_i}^q$ untuk setiap $i$ --- dan setelah
normalisasinya dibatalkan, $(\abs{a_i}^p)$ sebanding dengan $(\abs{b_i}^q)$.

\textbf{7.} Kasus $p = q = 2$ tak lain Cauchy--Schwarz dengan
kasus kesamaan yang sama (yakni kesebandingan). Untuk ujungnya:
$\sum\abs{a_ib_i} \leq \bigl(\max_i\abs{b_i}\bigr)\sum_i\abs{a_i} =
\norm a_1\norm b_\infty$, langsung suku demi suku.

\textbf{8.} Untuk $p = 1$ ia tak lain ketaksamaan segitiga suku demi
suku. Untuk $p > 1$, dengan $q$ yang sekawan:
\[
\norm{a+b}_p^p = \sum_i\abs{a_i + b_i}^p
\leq \sum_i\abs{a_i+b_i}^{p-1}\abs{a_i} +
\sum_i\abs{a_i+b_i}^{p-1}\abs{b_i},
\]
lalu Hölder pada tiap jumlahnya, dengan mencatat $(p - 1)q = p$:
\[
\sum_i\abs{a_i+b_i}^{p-1}\abs{a_i} \leq
\Bigl(\sum_i\abs{a_i+b_i}^{p}\Bigr)^{1/q}\norm a_p
= \norm{a + b}_p^{p/q}\,\norm a_p ,
\]
Demikian pula dengan $b$. Karenanya $\norm{a+b}_p^p \leq \norm{a +
b}_p^{p/q}\bigl(\norm a_p +
\norm b_p\bigr)$; lalu jika $a + b \neq 0$, bagilah dengan
$\norm{a+b}_p^{p/q}$ lalu pakai $p - \frac pq = 1$. Bersama
kehomogenan dan pemisahannya (yang jelas), $\norm\cdot_p$ merupakan
\omterm{def:b2:nvs:norm}{norma} pada $\R^n$.

\textbf{9.} Untuk $f, g$ yang \omterm{def:b2:metric:continuity}{kontinu} pada $\intcc ab$: Hölder memberi
\[
\int_a^b\abs{fg} \leq \Bigl(\int_a^b\abs
f^p\Bigr)^{1/p}\Bigl(\int_a^b\abs g^q\Bigr)^{1/q}
\]
lewat normalisasi yang sama ditambah Young titik demi titik yang
diintegralkan; sedangkan Minkowski $\norm{f + g}_p \leq \norm f_p +
\norm g_p$ diperoleh lewat pemecahan yang sama, dengan Hölder pada tiap
potongannya. Pemisahan \omterm{def:b2:nvs:norm}{normanya} memakai kepositifan tegas: sebab
$\abs f^p$ yang \omterm{def:b2:metric:continuity}{kontinu} dan berintegral nol pastilah nol secara identik
(jilid Tahun ke-1).

\textbf{10.} \emph{Kemonotonannya:} kita boleh mengandaikan $\norm a_p =
1$; maka tiap $\abs{a_i} \leq 1$, jadi $\abs{a_i}^q \leq \abs{a_i}^p$
sehingga $\norm a_q^q \leq 1$: yakni $\norm a_q \leq 1 = \norm a_p$.
Kesamaannya menuntut $\abs{a_i}^q = \abs{a_i}^p$ untuk setiap $i$,
yakni tiap $\abs{a_i} \in \{0, 1\}$; dan dengan $\sum\abs{a_i}^p = 1$
tersisa tepat satu koordinat bermodulus $1$: jadi kesamaannya bila dan
hanya bila $a$ punya paling banyak satu koordinat tak nol.
\emph{Limitnya:} $\norm a_\infty \leq \norm a_p \leq n^{1/p}\norm
a_\infty$, dan $n^{1/p} \to 1$.
\emph{Perbandingan sebaliknya:} Hölder dengan eksponen $\frac
qp$ beserta sekawannya $\frac{q}{q-p}$, yang diterapkan pada
$\abs{a_i}^p\cdot 1$, memberi
\[
\norm a_p^p = \sum_i\abs{a_i}^p\cdot 1 \leq
\Bigl(\sum_i\abs{a_i}^{q}\Bigr)^{p/q}\,n^{1 - p/q}
= \norm a_q^{p}\; n^{1-p/q},
\]
sehingga $\norm a_p \leq n^{\frac1p - \frac1q}\norm a_q$, dengan
kesamaan bila dan hanya bila semua $\abs{a_i}$ sama (yakni kasus
kesamaan Hölder terhadap vektor konstan).

\textbf{11.} Tulis $\abs{a_i}^r = \abs{a_i}^{\theta
r}\,\abs{a_i}^{(1-\theta)r}$ lalu terapkan Hölder dengan
eksponen sekawan $\frac{p}{\theta r}$ dan
$\frac{q}{(1-\theta)r}$ (yang sekawan justru karena $\frac{\theta
r}p + \frac{(1-\theta)r}q = 1$):
\[
\norm a_r^r = \sum_i \abs{a_i}^{\theta r}\abs{a_i}^{(1-\theta)r}
\leq \Bigl(\sum_i\abs{a_i}^{p}\Bigr)^{\theta r/p}
\Bigl(\sum_i\abs{a_i}^{q}\Bigr)^{(1-\theta)r/q}
= \norm a_p^{\theta r}\,\norm a_q^{(1-\theta)r} .
\]
Lalu ambil akar ke-$r$: jadi norma-$p$ bersifat log-cembung pada $\frac1p$.

\textbf{12.} \emph{Keduanya negatif:} jika $p < q < 0$ maka $0 < -q <
-p$, dan $M_{-q}(y) \leq M_{-p}(y)$ untuk eksponen positifnya
(yakni kasus di dalam pelajaran, \cref{ex:b2:realfun:powermeans}) yang
diterapkan pada $y = (1/x_i)$; lalu membalik kesamaan $M_p(x) =
M_{-p}(1/x)^{-1}$ membalik ketaksamaannya menjadi $M_p(x) \leq M_q(x)$.
\emph{Jembatannya:} untuk $q > 0$, kecekungan $\ln$ memberi $\ln M_q =
\frac1q\ln\bigl(\sum\lambda_ix_i^q\bigr) \geq
\frac1q\sum\lambda_i\ln x_i^q = \ln M_0$; sedangkan untuk $p < 0$,
kecekungan yang sama memberi $\ln\bigl(\sum\lambda_ix_i^p\bigr) \geq
p\sum\lambda_i\ln x_i$, lalu membaginya dengan $p < 0$ membaliknya: $\ln
M_p \leq \ln M_0$. Karenanya $M_p \leq M_0 \leq M_q$ setiap kali $p < 0
< q$: jadi bersama kedua kasus setandanya, $M$ naik pada seluruh
$\R^*$ (dan menembus $0$).

\textbf{13.} Misalkan $x_{\max} = \max x_i$, yang tercapai di $i^*$.
Untuk $p > 0$:
\[
\lambda_{i^*}^{1/p}\,x_{\max} \leq M_p \leq x_{\max},
\]
dan $\lambda_{i^*}^{1/p} \to 1$: jadi $M_p \to x_{\max}$. Untuk $p \to
-\infty$: $M_p(x) = M_{-p}(1/x)^{-1} \to
\bigl(\max_i\frac1{x_i}\bigr)^{-1} = \min_ix_i$.

\textbf{14.} Dengan $\lambda_i = \frac1n$, rantai $M_{-\infty}
\leq M_{-1} \leq M_0 \leq M_1 \leq M_2 \leq M_{+\infty}$ berbunyi
\[
\min \leq \frac{n}{\sum\frac1{a_i}} \leq \Bigl(\prod
a_i\Bigr)^{1/n} \leq \frac{\sum a_i}{n} \leq
\sqrt{\frac{\sum a_i^2}{n}} \leq \max .
\]
AM--HM (yakni $M_{-1} \leq M_1$) tertata ulang langsung menjadi
$\bigl(\sum a_i\bigr)\bigl(\sum\frac1{a_i}\bigr) \geq n^2$.

\textbf{15.} Dengan bobot sama, $M_p(x) =
\bigl(\frac1n\sum\abs{x_i}^p\bigr)^{1/p} = n^{-1/p}\norm x_p$.
Ketika $p$ membesar, $\norm x_p$ turun (pertanyaan 10) tetapi
penormalnya $n^{-1/p}$ naik lebih cepat, sehingga hasil kalinya
naik (pertanyaan 12): rata-rata itu merata-ratakan, \omterm{def:b2:nvs:norm}{norma} itu
menghimpun, dan faktor $n^{-1/p}$ persis menjadi kurs tukar antara
kedua kesepakatan pembukuan itu.

\textbf{16.} Tiap mata rantainya merupakan perwujudan Jensen yang tegas
(pertanyaan 2) dengan fungsi cembung atau cekung tegas pada
pertanyaan 3 (yakni $t^{q/p}$ dan $\ln$), jadi kesamaan pada mata rantai
mana pun memaksa semua $x_i$ sama; dan $\min = M_p$ atau $M_p = \max$
demikian pula memaksa semua nilainya sama dengan ekstremum bersamanya.
Jadi rantainya tegas begitu ada dua $x_i$ yang berbeda.

\textbf{17.} Terapkan Young (pertanyaan 4) pada pasangan
$\varepsilon^{1/p}a$ dan $\varepsilon^{-1/p}b$:
\[
ab = (\varepsilon^{1/p}a)(\varepsilon^{-1/p}b)
\leq \varepsilon\,\frac{a^p}p +
\varepsilon^{-q/p}\,\frac{b^q}q .
\]
Untuk $p = q = 2$, dengan mengganti $\varepsilon$ oleh $2\varepsilon$:
$ab \leq \varepsilon a^2 + \frac{b^2}{4\varepsilon}$ --- itulah
ketaksamaan penyerapan: sebuah hasil kali ditukar dengan kelipatan kecil
satu kuadrat ditambah kelipatan besar kuadrat lainnya.

\textbf{18.} Lewat teleskop:
\[
\prod_{k=1}^{n}c_k = \frac{\prod_{k=1}^n(k+1)^k}
{\prod_{k=1}^{n}k^{k-1}}
= \frac{2^1\,3^2\cdots(n+1)^n}{1^0\,2^1\cdots n^{n-1}}
= (n+1)^n,
\]
sebab tiap faktor $(k+1)^k$ pada pembilangnya tercoret dengan
suku berikutnya pada penyebutnya. Lalu AM--GM pada $n$ bilangan $c_ka_k$:
\[
(a_1\cdots a_n)^{1/n} =
\frac{\bigl(\prod_k c_ka_k\bigr)^{1/n}}{(n+1)}
\leq \frac{1}{n+1}\cdot\frac1n\sum_{k=1}^{n}c_ka_k .
\]

\textbf{19.} Setelah dijumlahkan atas $n$ lalu kedua penjumlahannya
ditukar (sebab semua sukunya positif: \cref{thm:b2:series:fubini}):
\[
\sum_{n\geq1}(a_1\cdots a_n)^{1/n}
\leq \sum_{n\geq1}\frac{1}{n(n+1)}\sum_{k=1}^{n}c_ka_k
= \sum_{k\geq1}c_ka_k\sum_{n\geq k}\frac1{n(n+1)}
= \sum_{k\geq1}\frac{c_ka_k}{k},
\]
dengan memakai teleskop $\sum_{n\geq k}\bigl(\frac1n -
\frac1{n+1}\bigr) = \frac1k$. Akhirnya $\frac{c_k}k =
\frac{(k+1)^k}{k^k} = \bigl(1 + \frac1k\bigr)^k < \eu$
(sebab barisan naik yang berlimit $\eu$, jilid Tahun ke-1):
\[
\sum_{n\geq1}(a_1\cdots a_n)^{1/n} \leq
\eu\sum_{k\geq1}a_k :
\]
itulah ketaksamaan Carleman. (Konstanta $\eu$ bersifat optimal, walaupun
kita tidak membuktikannya.)

\textbf{20.} Cauchy--Schwarz (pertanyaan 9 dengan $p = q = 2$) yang
diterapkan pada $\sqrt f$ dan $\frac1{\sqrt f}$ memberi
\[
1 = \Bigl(\int_0^1\sqrt f\cdot\frac{1}{\sqrt f}\Bigr)^{2}
\leq \Bigl(\int_0^1 f\Bigr)\Bigl(\int_0^1\frac1f\Bigr).
\]
Kesamaannya bila dan hanya bila $\sqrt f$ dan $\frac1{\sqrt f}$
sebanding, yakni $f^2$ konstan, yakni $f$ konstan (sebab $f > 0$ \omterm{def:b2:metric:continuity}{kontinu}).

\textbf{21.} Misalkan $1 < p < \infty$, $\norm a_p = \norm b_p = 1$,
$a \neq b$, dan andaikan $\bigl\Vert\frac{a+b}2\bigr\Vert_p = 1$,
yakni Minkowski menjadi kesamaan bagi $a, b$. Dengan merunut
bukti pertanyaan 8, kesamaannya memaksa kesamaan pada kedua penerapan
Hölder sekaligus pada ketaksamaan segitiga suku demi sukunya: jadi
$(\abs{a_i}^p)$ dan $(\abs{b_i}^p)$ keduanya sebanding dengan
$(\abs{a_i + b_i}^p)$, dan $a_i, b_i$ bertanda sama --- sehingga $b = ta$
untuk suatu $t \geq 0$, lalu $\norm b_p = \norm a_p$ memberi $t = 1$,
yakni $b = a$: kontradiksi. Karenanya permukaan bola-$p$ tak memuat
titik tengah dua titik permukaannya yang berbeda: jadi tanpa ruas.
Untuk $p = \infty$ di $\R^2$: semua $(1, t)$ dengan $\abs t \leq 1$
terletak pada permukaan bola satuannya --- yakni sebuah tepi datar;
sedangkan untuk $p = 1$: ruas $(t, 1 - t)$, $t \in \intcc01$, demikian pula.

\textbf{22.} Untuk $a = 0$ kedua ruasnya nol. Selain itu Hölder
membatasi setiap $\sum a_ib_i$ oleh $\norm a_p\norm b_q \leq \norm
a_p$. Ketercapaiannya: ambil
\[
b_i = \frac{\operatorname{sign}(a_i)\,\abs{a_i}^{p-1}}
{\norm a_p^{p/q}} :
\qquad
\norm b_q^q = \frac{\sum_i\abs{a_i}^{(p-1)q}}{\norm a_p^{p}}
= \frac{\norm a_p^p}{\norm a_p^p} = 1,
\quad
\sum_ia_ib_i = \frac{\norm a_p^p}{\norm a_p^{p/q}} =
\norm a_p ,
\]
dengan memakai $(p-1)q = p$ dan $p - \frac pq = 1$. Jadi supremumnya
adalah maksimum, sama dengan $\norm a_p$: sehingga tiap norma-$p$ adalah
\omterm{def:b2:nvs:norm}{norma} dual sekawannya --- yakni benih dualitas $L^p$--$L^q$.

\textbf{23.} Di sini $\E[X^r] = \sum_i\lambda_ix_i^r$, jadi
$\E[X^r]^{1/r} = M_r(x; \lambda)$, yang naik terhadap $r$ menurut
pertanyaan 12 (dan menembus $r \to 0, \pm\infty$ menurut pertanyaan
12--13): itulah ketaksamaan momen Lyapunov, yang murni pernyataan
tentang rata-rata pangkat berbobot. Ia kembali untuk peubah acak yang
sejati pada \cref{ch:b2:randomvar}.

\textbf{24.} (i) Rata-rata pangkat $M_1 \leq M_3$ dengan bobot sama
memberi $\frac{a+b+c}3 \leq \bigl(\frac{a^3+b^3+c^3}3\bigr)^{1/3}$;
lalu kubikkan dan kalikan dengan $3$: $a^3 + b^3 + c^3 \geq
\frac{(a+b+c)^3}9$. (ii) Cauchy--Schwarz terhadap vektor
konstan: $\sum_i\sqrt{x_i}\cdot1 \leq
\bigl(\sum_ix_i\bigr)^{1/2}n^{1/2}$; lalu kuadratkan.

\textbf{25.} Definisi tali busurnya menghasilkan lema kemiringan lewat
satu penataan ulang aljabar; lalu kemiringan yang diapit di sebuah titik
menghasilkan turunan sepihak dan garis penyangga, yang rata-rata
berbobotnya adalah Jensen. Diterapkan pada $-\ln$, Jensen menjadi Young,
yang setelah dijumlahkan terhadap vektor ternormalkan menjadi Hölder,
yang setelah dipecah lalu diserap ulang menjadi Minkowski --- dan
lahirlah norma-$p$ pada \cref{ch:b2:nvs}, \omterm{def:b2:metric:complete}{lengkap} dengan dualitasnya
(pertanyaan 22) dan geometrinya (pertanyaan 21). Jensen yang diterapkan
sepanjang skala pangkat merantai semua rata-rata dari $\min$ ke $\max$
(pertanyaan 12--14), yang dibaca pada peubah acak menjadi
ketaksamaan momen (pertanyaan 23). Dan AM--GM, setelah dibobot oleh satu
siasat teleskop, menghasilkan batas Carleman beserta konstantanya $\eu$
yang tak tereduksi (pertanyaan 18--19). Puncaknya:
Hölder--Minkowski, dan Carleman. Tujuannya: ruang $L^p$
pada jilid Tahun ke-3, yang aksioma pendirinya persis
pertanyaan 6 dan 8 dengan integral menggantikan jumlahnya.

\end{solution}
